
\documentclass{article}
\usepackage{colortbl}
\usepackage{makecell}
\usepackage{multirow}
\usepackage{supertabular}

\begin{document}

\newcounter{utterance}

\centering \large Interaction Transcript for game `wordle\_withclue\_noexp', experiment `medium\_frequency\_words\_clue\_no\_critic', episode 8 with reflexion{-}Qwen3.5{-}9B.
\vspace{24pt}

{ \footnotesize  \setcounter{utterance}{1}
\setlength{\tabcolsep}{0pt}
\begin{supertabular}{c@{$\;$}|p{.15\linewidth}@{}p{.15\linewidth}p{.15\linewidth}p{.15\linewidth}p{.15\linewidth}p{.15\linewidth}}
    \# & \multicolumn{2}{c}{Player} && \multicolumn{2}{c}{Game Master} \\
    \hline

    \theutterance \stepcounter{utterance}  
    & & & \multicolumn{4}{p{0.6\linewidth}}{
        \cellcolor[rgb]{0.9,0.9,0.9}{
            \makecell[{{p{\linewidth}}}]{
                \texttt{\tiny{[P1$\langle$GM]}}
                \texttt{You have attempted this same task before and did not succeed. These are the plans you wrote after those attempts.} \\
\texttt{Trial \#1: New plan: I failed because I guessed "make", which is only four letters long, violating the rule that the target word must be five lowercase letters. Next time, I will ensure my guess is exactly five letters long and matches the clue "make current" (e.g., "maker" or "current"), while strictly adhering to the format "guess: {[}word{]}" without adding any explanation tags or extra text.} \\
\\ 
\texttt{Trial \#2: New plan: I failed because I included the clue word "current" as a guess, which is only six letters long, violating the rule that the target word must be exactly five lowercase letters. Additionally, I violated the strict formatting rule by adding an explanation tag ("explanation:") which was not requested. Next time, I will strictly adhere to the format "guess: {[}word{]}" with no extra text, and I will ensure my guess is a valid five{-}letter English word that incorporates the letters from the clue "make current" (m, a, k, e, r, c, u, n, t) while excluding any red letters (m, a, k) and utilizing the green (e) and yellow (r) letters correctly. A valid five{-}letter guess could be "cater" or "crane" or "carte".} \\
\\ 
\texttt{Trial \#3: New plan: I failed because I got stuck in an infinite loop by repeatedly guessing "crane" after receiving the same feedback, wasting all my attempts without exploring other possibilities. The feedback `c<red> r<yellow> a<red> n<yellow> e<yellow>` indicates that 'e' is correct and in the right position, while 'r', 'n' are correct but in the wrong positions, and 'c', 'a' are incorrect. I should have immediately used this information to construct a new five{-}letter word that keeps 'e' at the end, places 'r' and 'n' in different slots, and excludes 'c' and 'a'. For example, I could have tried "nerve" (keeping 'e' at the end, moving 'r' and 'n') or "never" or "renew". Next time, upon receiving feedback, I will immediately generate a distinct new guess that incorporates green letters in their fixed positions, moves yellow letters to unoccupied slots, and replaces red letters, ensuring I do not repeat the same word.} \\
\\ 
\\ 
\texttt{You are a language wizard who likes to guess words by using the given rules.} \\
\\ 
\texttt{Welcome to Wordle! You have six attempts to guess the target word, a valid English word of five lowercase letters (a{-}z). Please use the tag "guess:" to provide your guess.} \\
\\ 
\texttt{To help you make an informed guess, you will receive a clue for the word, such as} \\
\texttt{clue: snowy white.} \\
\\ 
\texttt{Here is an example guess based on the clue: } \\
\texttt{guess: apple} \\
\\ 
\texttt{After each guess, your answer will be validated, and you will receive feedback indicating which letters are correct (green), which letters are correct but in the wrong position (yellow), and which letters are incorrect (red). This feedback can be useful in determining which letters to include or exclude in your next guess.} \\
\\ 
\texttt{For example, the feedback for "apple" might be:} \\
\texttt{guess\_feedback: a<yellow> p<yellow> p<green> l<yellow> e<red>} \\
\\ 
\texttt{Do not generate any other text.} \\
\\ 
\texttt{Let's begin with your first guess.} \\
\\ 
\texttt{clue: make current} \\
            }
        }
    }
     \\ \\

    \theutterance \stepcounter{utterance}  
    & \multicolumn{4}{p{0.6\linewidth}}{
        \cellcolor[rgb]{0.561,0.737,0.561}{
            \makecell[{{p{\linewidth}}}]{
                \texttt{\tiny{[P1$\rangle$GM]}}
                \texttt{guess: current} \\
            }
        }
    }
    & & \\ \\

    \theutterance \stepcounter{utterance}  
    & & & \multicolumn{2}{p{0.3\linewidth}}{
        \cellcolor[rgb]{0.9,0.9,0.9}{
            \makecell[{{p{\linewidth}}}]{
                \texttt{\tiny{[GM$|$GM]}}
                \texttt{Error: The length of the guessed word is not 5.} \\
            }
        }
    }
    & & \\ \\

    \theutterance \stepcounter{utterance}  
    & & & \multicolumn{2}{p{0.3\linewidth}}{
        \cellcolor[rgb]{0.9,0.9,0.9}{
            \makecell[{{p{\linewidth}}}]{
                \texttt{\tiny{[GM$|$GM]}}
                \texttt{game\_result = ABORT} \\
            }
        }
    }
    & & \\ \\

    \theutterance \stepcounter{utterance}  
    & & \multicolumn{4}{p{0.6\linewidth}}{
        \cellcolor[rgb]{0.95,0.95,0.95}{
            \makecell[{{p{\linewidth}}}]{
                \texttt{\tiny{[P1$\leftrightarrow$P]}}
                \texttt{Post{-}episode Reflexion (not sent to GM):} \\
\texttt{New plan: I failed because I guessed "current", which is six letters long, violating the strict rule that the target word must be exactly five lowercase letters. Additionally, I ignored the specific formatting instruction to use only the tag "guess:" without any explanatory text. Next time, I will strictly adhere to the format "guess: {[}word{]}" with no extra text, ensuring my guess is a valid five{-}letter English word derived from the clue "make current" (e.g., "maker", "cater", "crane", or "carte") that fits the length constraint.} \\
            }
        }
    }
    & \\ \\

\end{supertabular}
}

\end{document}
